\subsection{Cortex-A57 处理器}

现代 Arm Cortex 与 Neoverse 产品已经远超 Cortex-A57 所处的代际，
但 A57 是较早、资料完整的乱序 Armv8-A 核，适合用来理解经典 AArch64 实现。
本节讨论的是 Cortex-A57 的具体微架构，而不是所有 Armv8-A 处理器的通用属性。

Cortex-A57 面向移动和企业计算应用，可与 Cortex-A53 组成早期 Arm big.LITTLE 系统，
在性能与功耗之间进行任务迁移或异构调度。

Cortex‑A57 处理器具有与其他处理器的高速缓存一致性互操作性，包括用于 GPU 计算的 ARM Mali™ 系列图形处理单元 (GPU)，并为高性能企业应用程序提供可选的可靠性和可扩展性功能。
它提供了比 ARMv7 架构的 Cortex‑A15 处理器更高的性能，并具有更高的能效水平。
可选密码扩展可显著加速部分算法，但实际收益依赖具体实现与工作负载。

\begin{figure}[htbp]
  \centering
  \resizebox{0.72\textwidth}{!}{%
    \begin{tikzpicture}[
      font=\sffamily\small,
      block/.style={draw=booknavy, rounded corners=2mm, fill=bookpaper,
        align=center, minimum height=0.72cm},
      unit/.style={block, fill=white, minimum height=1.15cm},
      cluster/.style={draw=booknavy, rounded corners=2mm, fill=bookline,
        minimum width=8.4cm, minimum height=5.0cm}
    ]
      \fill[bookblue!55, rounded corners=3mm] (0,0) rectangle (10.2,10.9);
      \node[block, minimum width=9.5cm] at (5.1,10.35)
        {CoreSight multicore debug and trace};
      \node[block, minimum width=9.5cm] at (5.1,9.45)
        {Generic Interrupt Controller};

      \node[cluster] at (5.75,6.55) {};
      \node[cluster] at (5.4,6.25) {};
      \node[cluster] at (5.05,5.95) {};
      \node[cluster, fill=bookpaper!80!bookline] at (4.7,5.65) {};
      \node[font=\sffamily\scriptsize\bfseries, anchor=east] at (8.72,4.05)
        {Cores 0--3};

      \node[unit, minimum width=4.6cm, minimum height=2.75cm] at (3.25,6.75)
        {Cortex-A57\\out-of-order core};
      \node[unit, minimum width=2.25cm, minimum height=1.3cm] at (6.9,7.5)
        {Advanced SIMD\\and crypto};
      \node[unit, minimum width=2.25cm, minimum height=1.05cm] at (6.9,6.0)
        {Floating-point\\unit};
      \node[unit, minimum width=2.25cm, minimum height=1.05cm] at (6.9,4.65)
        {Memory\\protection unit};
      \node[unit, minimum width=2.15cm, minimum height=1.05cm] at (1.95,4.65)
        {L1 instruction\\cache};
      \node[unit, minimum width=2.15cm, minimum height=1.05cm] at (4.3,4.65)
        {L1 data cache\\with ECC};
      \node[unit, minimum width=4.5cm, minimum height=0.75cm] at (3.25,3.55)
        {Performance Monitor Unit};

      \node[block, minimum width=1.9cm] at (1.45,2.25) {SCU};
      \node[block, minimum width=1.9cm] at (3.65,2.25) {ACP};
      \node[block, minimum width=5.0cm] at (7.3,2.25)
        {Integrated shared L2 cache with ECC};
      \node[block, minimum width=9.5cm, minimum height=0.8cm] at (5.1,0.85)
        {AMBA 4 ACE coherent interface};
    \end{tikzpicture}%
  }
  \caption{Cortex-A57 集群结构（教学简图）}
  \label{fig:cortex-a57}
\end{figure}

Cortex-A57 实现 Armv8.0-A，并支持 AArch64 和 AArch32 执行状态。
它支持多核操作，在单个集群中具有一到四核多处理器。
该核采用 AMBA 4 ACE 一致性接口；系统级互连是否使用其他协议取决于具体 SoC。
CoreSight 组件可用于调试和跟踪。

Cortex‑A57 处理器具有以下特性：

\begin{itemize}
  \item 乱序，多于 15 级流水线。
  \item 省电功能包括路径预测、标记缩减和缓存查找抑制。
  \item 通过重复执行资源增加峰值指令吞吐量。
    具有本地化解码、3 个宽解码带宽的功率优化指令解码。
  \item 性能优化的 L2 缓存设计使集群中的多个核心可以同时访问 L2。
\end{itemize}
